\chapter{Slopes and folds}\label{ch:slopes-folds}
\begin{goals}
\textbf{In this chapter} you test a rock slope for planar, wedge and toppling failure, and describe a fold
from its bedding: axis, interlimb angle, axial plane, and Fleuty's names for them.
\end{goals}

\section{Kinematic analysis of a slope}\label{sec:kinematic}\index{kinematic analysis}

\begin{center}\small\ttfamily
kinematic joints 300/65 --phi 30 --lateral 20
\end{center}

\noindent (\menu{Analysis>Slope stability (kinematic)…} asks for the layer, the face, $\phi$ and the
limit) tests the planes of \code{joints} against a slope face dipping 65° toward 300, with a friction
angle of 30° and a lateral limit of ±20° (the defaults are 30° and 20°). Each mode is tested on the
attitudes themselves \parencite{markland1972,hoek1981,goodman1976}, so the result doesn't depend on the
projection:

\begin{description}
\item[planar] a plane that dips out of the face (its dip direction within the lateral limit of the
face's), less steeply than the face (so it daylights) and more steeply than $\phi$.
\item[wedge] two planes whose line of intersection plunges out of the face, less steeply than the face
in that direction, and more steeply than $\phi$. Every pair is tested; past 50\,000 pairs wuffle asks you to
use the set means or a selection instead.
\item[flexural toppling] a plane that dips into the slope (within the lateral limit of the face's
opposite), steep enough that its pole plunges less than the face's dip less $\phi$.
\end{description}

The answer counts each mode, and every plane gets a \col{kinematic} column naming the modes it takes part
in (\code{planar wedge}), so \code{select where "kinematic \textasciitilde{} 'wedge'"} picks them. The face
and the friction cone ($90 - \phi$ about the vertical) are drawn on the net as datasets of their own. The
whole analysis is one undo step.

\section{Folds}\label{sec:fold}\index{folds}

\code{fold} (\menu{Analysis>Fold…}) describes a fold from bedding:

\begin{description}
\item[The axis] the pole of the best-fit great circle through the bedding poles (the $\pi$-girdle), with
its cylindricity: how well the poles sit on one great circle.
\item[The limbs] the two sets the poles fall into.
\item[The interlimb angle] measured across the hinge, and Fleuty's word for it: gentle (180–120°), open
(120–70°), close (70–30°), tight, isoclinal \parencite{fleuty1964}.
\item[The axial plane] the plane that holds the axis and bisects the limbs, and Fleuty's names for the
fold's attitude (upright, inclined, recumbent) and its plunge (horizontal to vertical).
\end{description}

\begin{caution}
Two limbs have two bisecting planes, both holding the axis. Which one is the axial plane is a fact about
the fold, not about the orientations. wuffle takes the steeper, which is right for every fold but a
recumbent one, and says so. For a recumbent fold, \code{fold --recumbent} takes the other, and the
interlimb angle changes with it.
\end{caution}

\noindent \code{fold --keep} draws the axis and the axial plane as datasets. When the poles sit poorly on
one great circle, the answer suggests \code{cone}: the poles of a conical fold lie on a small circle, and
\code{cone} (\menu{Analysis>Best-fit cone…}) finds the one that fits best (its axis, half-angle and the
scatter about it), drawn with \code{--keep}.
